\chapter{Modelo de planificacion}
\label{anexo3:cap:modelo}
Para la planificación del proyecto se ha utilizado un modelo ágil, específicamente Scrum, una
metodología ágil que se basa en la entrega incremental y la colaboración continua entre
el equipo de desarrollo y los interesados. Se ha optado por un enfoque con una gestión ágil en lugar
de una gestión predictiva por varias razones. La principal es la falta de experiencia previa a
la hora de realizar estimaciones de esfuerzo y tamaño, especialmente estimaciones muy a largo
plazo, principalmente debido a la falta de experiencia en el desarrollo de software complejo:
no puedo estimar con precisión el esfuerzo necesario para compeltar tareas que aun no sé
qué tan complejas van a ser.

El desarrollo ágil permite una mayor flexibilidad y adaptabilidad a los cambios, lo que es
especialmente útil en proyectos donde los requisitos pueden evolucionar a lo largo del tiempo
\parencites{AgileManifesto}[p.~16]{abrahamsson2017agile}. La planificación ágil resuelve justo
el problema de la incapacidad de estimar el esfuerzo necesario para completar la totalidad del
proyecto al inicio, limitándose a estimar el esfuerzo de las tareas más inmediatas
\parencite[p.~2]{robert2003agile}.

La flexibilidad del modelo ágil surge del hecho de que el proyecto se divide en iteraciones muy
cortas o sprints, lo que permite que el equipo de desarrollo se enfoque en un conjunto reducido
de tareas a la vez \parencite{AgileManifesto}. Esto no solo facilita la planificación y el
seguimiento del progreso, sino que también permite una retroalimentación continua y la posibilidad
de ajustar el enfoque según las necesidades cambiantes del proyecto. Además, puesto que estos sprints
son cortos y frecuentes, el equipo puede internalizar rápidamente la cantidad de trabajo que se puede
realizar en un período de tiempo determinado, lo que mejora la capacidad de estimación a medida que
el proyecto avanza \parencite[p.~3]{robert2003agile}.

\diagrama{
    nombre = {Planificación ágil (Imagen de Planbox, Creative Commons)},
    archivo = {anexo3/imgs/agile.png},
    etiqueta = {anexo3:fig:agile},
}

\section{Scrum}
\label{anexo3:sec:scrum}
Scrum es un marco de trábajo ligero y ágil que se utiliza para gestionar proyectos de
desarrollo de software \parencite{schwaber2020scrum}. Scrum se basa en la idea de que el
conocimiento proviene de la experiencia y que se deben tomar decisiones basadas en hechos
observables y no en suposiciones. Además, Scrum fomenta la autoorganización e intenta reducir
el desperdicio de tiempo y recursos al enfocarse en la entrega continua de valor al cliente.

\subsection{Sprints}
\label{anexo3:subsec:sprints}
Scrum organiza el trabajo en ciclos cortos llamados sprints, que suelen durar entre una y cuatro
semanas. Cada sprint se debe considerar como un mini-proyecto que tiene un objetivo claro y
un conjunto de tareas a completar. Aumentar o disminuer la duración de los sprints afecta de
manera proporcional a la cantidad de trabajo realizado, pero también al riesgo de que dicho
trabajo no sea de calidad o no cumpla con los requisitos del cliente \parencite{schwaber2020scrum}.

Un sprint comienza con una planificación, donde el equipo selecciona un conjunto de tareas del
backlog del producto que se compromete a completar durante el sprint. Durante el sprint, el
equipo trabaja en las tareas seleccionadas y realiza reuniones diarias para revisar el progreso
y ajustar el enfoque si es necesario. Al final del sprint, ocurren el \textit{Sprint Review} y el
\textit{Sprint Retrospective}. En el \textit{Sprint Review}, el equipo presenta el trabajo realizado
y recibe retroalimentación del cliente y otros interesados. En el \textit{Sprint Retrospective}, el
equipo reflexiona sobre el sprint y discute qué funcionó bien, qué no funcionó y cómo se puede
mejorar en el próximo sprint \parencite{schwaber2020scrum}.

\subsection{Artefactos}
\label{anexo3:subsec:artefactos}
Scrum define tres artefactos principales: el Product Backlog, el Sprint Backlog y el Incremento.
El Product Backlog es una lista priorizada de las tareas necesarias para mejorar el estado del
producto, que se actualiza continuamente a medida que se identifican nuevas tareas o se transladan
tareas a los sprints. El Sprint Backlog es una lista de las tareas seleccionadas para un sprint
concreto, que se compromete a completar durante ese sprint. El Incremento es el resultado del
trabajo realizado durante un sprint, que debe ser un producto funcional y potencialmente
liberable al cliente \parencite{schwaber2020scrum}.

Para representar gráficamente el proceso del sprint se ha utilizado un panel similar al de Kanban
(ver Figura \ref{anexo3:fig:kanban}), doonde cada columna representa un estado del proceso
como puede ser ``Por hacer'', ``En proceso'' y ``Hecho''. Este panel permite visualizar el flujo de
trabajo y el progreso del equipo durante el sprint, facilitando la identificación de cuellos de botella
y la gestión del trabajo pendiente.

\diagrama{
    nombre = {Panel de Kanban (Imagen de Jeff.Iasovski, Creative Commons)},
    archivo = {anexo3/imgs/kanban.png},
    etiqueta = {anexo3:fig:kanban},
}

\subsection{Roles}
\label{anexo3:subsec:roles}
Scrum define tres roles principales: el Product Owner, el Scrum Master y el equipo de desarrollo.
El Product Owner es responsable de definir y priorizar los requisitos del producto, asegurando
que el equipo trabaje en las tareas más importantes para el cliente. El Scrum Master es el encargado
de facilitar el proceso de Scrum, eliminando obstáculos y asegurando que el equipo siga las gúias
establecidas. El equipo de desarrollo es el que realiza el trabajo técnico, desarrollando el software
y cumpliendo los requisitos del Product Owner \parencite{schwaber2020scrum}.

% \section{Kanban}
% \label{anexo3:sec:kanban}
% Kanban es un método ágil para gestionar el trabajo y mejorar la eficiencia en los procesos de
% desarrollo de software \parencite{LAGEJUNIOR201013}. Se basa en la visualización del flujo de trabajo
% en columnas (ver Figura \ref{anexo3:sec:kanban}), donde cada una representa un estado del proceso como
% puede ser ``Por hacer'', ``En proceso'' y ``Hecho''. Además, tiene varias diferencias con Scrum, como
% la ausencia de roles definidos o la falta de sprints.

% \section{Scrumban}
% \label{anexo3:sec:scrumban}
% Scrumban es un enfoque híbrido que combina elementos de Scrum y Kanban, diseñado para
% proporcionar la flexibilidad de Kanban con la estructura de Scrum \parencite{Stoica2016AnalyzingAD}.
% De la misma manera que Scrum, Scrumban utiliza sprints para organizar el trabajo en ciclos cortos
% (a poder ser más cortos que en Scrum), pero también incorpora la visualización del flujo de trabajo
% de Kanban para mostrar el estado de las tareas en cada sprint. Kanban proporciona un límite de trabajo
% en cada sección del flujo, que se cumple con el objetivo de evitar la sobrecarga del equipo. Esto
% se consigue, por ejemplo, limitando el número de tareas que pueden estar en la columna ``En proceso''
% o reduciendo la longitud de los sprints para disminuir la cantidad de trabajo total del ciclo
% \parencite{Scrumban}.
%
% En Scrumban, al igual que en Kanban, no existen roles definidos, haciendo que todos los miembros del
% equipo de desarrollo sean responsables de la planificación y ejecución del trabajo. Esto implica,
% entre otras cosas, que es el propio equipo el encargado de priorizar las tareas y decidir cuáles
% se deben realizar en cada sprint \parencite{Scrumban}.

